\section{Enforcing Non-Negativity}\label{sec:primaldual}

Section~\ref{sec:pd} reduces the long-only problem to a handful of SPD solves, leaving
$w \geq 0$ as the sole remaining constraint. It is enforced by an \emph{active-set}
strategy that separates the balance constraints from the non-negativity
constraint: the balance-only problem is solved over a dynamically adjusted \emph{active
set} of assets, and non-negativity is enforced iteratively without ever modifying the
inner solve.

In the \emph{primal step}, any asset with a negative weight is dropped from the active
set and the balance-constrained SPD reduction is re-solved on the reduced universe. Once
all active weights are non-negative, the \emph{dual step} checks the KKT gradient
condition for every excluded asset: global optimality requires every excluded asset $i$
to satisfy
\[
  \nabla_i f(w) \;\geq\; (B^\top\lambda)_i\;,
  \qquad
  \nabla_i f(w) = \tfrac{2}{T}(X^\top Xw)_i - \rho\hat\mu_i\;,
\]
where $\lambda \in \mathbb{R}^p$ collects the balance multipliers recovered analytically
from Section~\ref{sec:pd} (in the budget case, the single shadow price of full
investment). An excluded asset that violates this condition would decrease
portfolio variance if held at a small positive weight; it is re-added to the active set
and the loop repeats. The loop terminates when both primal and dual feasibility hold;
together with stationarity from the inner solve this certifies global optimality.

These two conditions are precisely the complementarity system for the long-only KKT
conditions: with the constraint multiplier $s_i = \nabla_i f(w) - (B^\top\lambda)_i$, an
optimum requires $w \geq 0$, $s \geq 0$, and $w_i s_i = 0$, and the asset toggles are the
\emph{principal pivots} of the associated linear complementarity problem
\cite{judice1994,kim2011,murty1988}. This loop, block principal pivoting on a
$P$-matrix LCP run as a fast path guarded by a least-index (Bland-style) anti-cycling
fallback, is developed and analysed in full generality in the companion
paper~\cite{schmelzer2026nncg}, for the non-negative quadratic
$\min_{x \geq 0} \tfrac{1}{2}x^\top A x - b^\top x$ with $A \succ 0$ and its
equality-augmented variant. Algorithm~\ref{alg:elim} is that loop instantiated on the
balance-constrained minimum-variance problem: let $N$ be the total number of primal and
dual violators at the current solve and $\bar N$ the fewest seen so far; a batch
exchange is taken whenever it makes progress ($N < \bar N$) or a fixed patience budget
has not been exhausted, and otherwise a single least-index pivot is taken.

\begin{algorithm}[ht]
\caption{Active-set outer loop (wraps any inner solver $f$)}\label{alg:elim}
\begin{algorithmic}[1]
\Require Inner solver $f$;\; tolerance $\varepsilon = 10^{-6}$;\; $\alpha > 0$
  \hfill\textit{(ensures $\hat\Sigma_{\mathrm{LW},\mathcal{A}}$ non-singular)};\; patience $p_{\max} \in \mathbb{N}$ \textit{(default $3$)}
\Ensure Weight vector $w \in \mathbb{R}^n$;\; outer step count $s$;\; total inner iteration count $k$ \textit{(equals $s$ for direct solvers using the $k_{\mathrm{step}}=1$ sentinel)}
\State $\mathrm{active} \leftarrow \{1,\ldots,n\}$;\; $s \leftarrow 0$;\; $k \leftarrow 0$;\; $\bar N \leftarrow n+1$;\; $p \leftarrow p_{\max}$
  \hfill\textit{($\bar N$: fewest infeasibilities seen; $p$: batch-step budget)}
\Loop
  \State $(w_{\mathrm{a}},\, k_{\mathrm{step}}) \leftarrow f(\mathrm{active})$;\; $s \leftarrow s + 1$;\; $k \leftarrow k + k_{\mathrm{step}}$
    \hfill\textit{($k_{\mathrm{step}} = 1$ sentinel for direct solvers; iterative count for CG)}
  \State $w_i \leftarrow w_{a,i}$ for $i \in \mathrm{active}$;\; $w_i \leftarrow 0$ for $i \notin \mathrm{active}$
  \State $g_i \leftarrow \tfrac{2}{T}(X^\top Xw)_i - \rho\hat\mu_i$;\;
         $\lambda \leftarrow$ balance multipliers from the inner solve (Section~\ref{sec:pd})
    \hfill\textit{(stationarity: $g_j \approx (B^\top\lambda)_j$ for $j \in \mathrm{active}$ to solver tolerance)}
  \State $\mathcal{D} \leftarrow \{i \in \mathrm{active} : w_i < -\varepsilon\}$;\;
         $\mathcal{V} \leftarrow \{i \notin \mathrm{active} : g_i < (B^\top\lambda)_i - \varepsilon\}$
    \hfill\textit{(primal and dual violators)}
  \State $\mathcal{H} \leftarrow \mathcal{D} \cup \mathcal{V}$;\; $N \leftarrow |\mathcal{H}|$
  \If{$N = 0$}
    \State \textbf{break} \hfill\textit{(KKT satisfied; globally optimal)}
  \ElsIf{$N < \bar N$ \textbf{ or } $p > 0$}
    \hfill\textit{(fast path: progress, or patience remains)}
    \State \textbf{if} $N < \bar N$ \textbf{ then } $\bar N \leftarrow N$;\; $p \leftarrow p_{\max}$ \textbf{ else } $p \leftarrow p - 1$
    \State $\mathrm{active} \leftarrow (\mathrm{active} \setminus \mathcal{D}) \cup \mathcal{V}$
      \hfill\textit{(batch exchange: drop all $\mathcal{D}$, add all $\mathcal{V}$)}
  \Else
    \hfill\textit{(anti-cycling fallback: $p$ exhausted)}
    \State $i^\star \leftarrow \min\{i : i \in \mathcal{H}\}$
      \hfill\textit{(Bland least-index rule)}
    \State \textbf{if} $i^\star \in \mathcal{D}$ \textbf{ then } $\mathrm{active} \leftarrow \mathrm{active} \setminus \{i^\star\}$
           \textbf{ else } $\mathrm{active} \leftarrow \mathrm{active} \cup \{i^\star\}$
      \hfill\textit{(single principal pivot)}
  \EndIf
\EndLoop
\State \Return $w$, $s$, $k$
\end{algorithmic}
\end{algorithm}

\paragraph{Balance multipliers}
The multipliers in the $\lambda$-computation step of Algorithm~\ref{alg:elim} come for
free: the Schur elimination of Section~\ref{sec:pd} produces $\lambda$ alongside the
weights of every inner solve. In the budget case the implementation instead recovers
the single shadow price directly from the gradients: in exact arithmetic the inner
solver minimises $f$ over $\{w_\mathcal{A} : \mathbf{1}^\top w_\mathcal{A} = 1\}$, whose
stationarity condition is
$\tfrac{2}{T}(X_\mathcal{A}^\top X_\mathcal{A}\,w_\mathcal{A})_j = \lambda$
for every $j \in \mathcal{A}$, so all active gradients equal $\lambda$.
In practice CG terminates at residual tolerance $\varepsilon = 10^{-5}$, so
the active gradients agree with $\lambda$ up to this tolerance; the median (used in the implementation for robustness; the mean gives the
same value in exact arithmetic) recovers $\lambda$ to the same tolerance.

\begin{theorem}[Finite convergence of Algorithm~\ref{alg:elim}]\label{thm:termination}
Let $f(w) = w^\top \hat\Sigma_{\mathrm{LW}}\,w - \rho\hat\mu^\top w$ with
$\hat\Sigma_{\mathrm{LW}} \succ 0$ (ensured by $\alpha > 0$).
Then for any patience budget $p_{\max} \in \mathbb{N}$, Algorithm~\ref{alg:elim}
terminates in finitely many outer iterations and returns the unique global minimizer
of $\min_{w}\,f(w)$ subject to $Bw = c,\;w \geq 0$, provided $B_{\mathcal{A}}$ retains
full row rank on every active set the loop visits (automatic for the budget case
$p = 1$; generic once $|\mathcal{A}| \geq p$).
No non-degeneracy assumption is required.
\end{theorem}

\begin{proof}
Because $\hat\Sigma_{\mathrm{LW}} = (1-\alpha)\hat\Sigma + \alpha T_0 \succ 0$, every
principal submatrix $\hat\Sigma_{\mathrm{LW},\mathcal{A}}$ is SPD, and with
$B_{\mathcal{A}}$ of full row rank the Schur elimination of Section~\ref{sec:pd} is
well posed at every outer step. The problem is then the equality-augmented
non-negative quadratic treated in \cite{schmelzer2026nncg}, and
Algorithm~\ref{alg:elim} is the loop analysed there. Theorem~4.1 of
\cite{schmelzer2026nncg} gives finite termination at the unique global minimiser; its
exit condition is precisely the $N = 0$ test of Algorithm~\ref{alg:elim}, which
certifies the long-only KKT conditions.
\end{proof}

The guarantee is unconditional (the least-index fallback removes the non-degeneracy
hypothesis that classical block-principal-pivoting arguments require) and it is
robust to the inexactness of the CG inner solves: once CG's residual tolerance is tied
to the KKT decision threshold $\varepsilon$, the inexact loop provably visits the same
active sets as the exact one \cite[Lemma~4.1]{schmelzer2026nncg}. The fallback is a
safety net that generic data never triggers: a batch step fails to make progress only
on a codimension-one, hence non-generic, configuration of $X$, though an adversarial
family of anti-correlated designs shows the guard is genuinely necessary
\cite[Section~6.2]{schmelzer2026nncg}. In practice the outer loop never reaches the
fallback: across both equity universes and all shrinkage levels reported here it
terminates in $s$ between $6$ and $18$ outer steps
(Tables~\ref{tab:sp500}--\ref{tab:ftse}), with the active set changing by only a few
assets per step and the patience budget $p$ never exhausted, so the batch fast path
alone certifies optimality and no cycling was observed in any run.

\paragraph{Warm-starting}
Two levels of warm-starting are available for Algorithm~\ref{alg:elim}.

\emph{Within-solve warm-starting} (between outer iterations of a single solve).
When the active set changes size, the previous CG solution lives in a different
subspace and cannot be reused directly. One can restrict the previous full
weight vector $w^{(k)} \in \mathbb{R}^n$ to the new active set
$\mathcal{A}^{(k+1)}$ and, in the budget case, renormalise:
\[
  x_0 = \frac{w^{(k)}_{\mathcal{A}^{(k+1)}}}
             {\mathbf{1}^\top w^{(k)}_{\mathcal{A}^{(k+1)}}}\,.
\]
This approximates $\hat\Sigma_{\mathcal{A}^{(k+1)}}^{-1}\mathbf{1}$ up to a
positive scalar, which is what CG seeks.
The benefit is limited within a single solve because the outer loop terminates in few outer steps in practice (6--18 on S\&P~500 data; see Table~\ref{tab:sp500}).

\emph{Cross-solve warm-starting} (across a sequence of related problems).
For an efficient-frontier sweep $\rho_1 < \rho_2 < \cdots < \rho_N$, consecutive
optimal portfolios differ by only a few assets entering or leaving the active
set.
Passing the previous problem's final $(\mathcal{A},\,w)$ pair as the
initial state for the next solve simultaneously reduces outer iterations (the
active set needs fewer primal and dual adjustments) and CG iterations (the
initial residual is small). When the optimal support does not move between
adjacent $\rho$ points, the warm-started loop terminates in a single outer step
and the inner iteration count scales with the logarithm of the parameter step
\cite[Proposition~4.2]{schmelzer2026nncg}.
The mechanism is the same for a direct solver
such as KKT, but KKT profits only from the warm active set: direct
factorisation has no analogue of an initial guess.

The cross-solve warm start is implemented in \texttt{solve\_cg\_warm}
(warm active set $+$ CG initial guess) and \texttt{solve\_kkt\_warm}
(warm active set only).  Section~\ref{sec:results} benchmarks \texttt{solve\_cg\_warm}
against cold-started CG on a 50-point frontier sweep.

\medskip
A loop-free alternative (projected gradient, which enforces the constraints at
every step by projecting onto the simplex) is analysed in the companion
paper~\cite{schmelzer2026nncg}: it is not a Krylov method, forgoes the subspace
structure, and converges at $O(\kappa)$ against the $O(\sqrt{\kappa})$ of the CG inner
solver at the same $O(nT)$ per-step cost. Section~\ref{sec:results} benchmarks it as a
matrix-free first-order baseline.
Section~\ref{sec:solvers} develops the concrete implementations.
